\section{Method}



We formulate test-time alignment as feedback control over the observable
attributes of a frozen large language model (LLM). Our approach learns a
finite-memory Koopman predictor of prompt--response dynamics and uses it
to evaluate interventions before querying the LLM. The predictor maps
recent attribute observations into a learned observable space, propagates
them through affine dynamics, and decodes the resulting attribute
trajectories. At inference time, the controller selects an intervention
according to its predicted alignment cost, executes it, and replans from
the newly observed response. The LLM parameters remain fixed throughout;
only the external dynamics model is learned offline.

\begin{figure}[ht]
    \centering
    \includegraphics[width=0.95\linewidth]{figs/method.png}
    \caption{Caption}
    \label{fig:placeholder}
\end{figure}

\subsection{Problem Formulation}
\label{sec:problem}

Let $x$ denote a task input, $p_{\theta}$ a frozen LLM, and
$r\in\mathbb{R}$ the desired value of an output attribute. We consider
one attribute at a time, such as sentence length or sentiment polarity.
At interaction step $t$, the controller observes the dialogue history
$\mathcal{H}_t$ and chooses a textual intervention $u_t$. The next
response and its measured attribute are
\begin{equation}
    o_{t+1}\sim
    p_{\theta}(\,\cdot\mid x,\mathcal{H}_t,u_t),
    \qquad
    y_{t+1}=V(o_{t+1}),
    \label{eq:interaction}
\end{equation}
where $V$ is a task-specific verifier. We distinguish the generated
text $o_t$ from its numerical attribute $y_t$. A time step corresponds
to an entire prompt--response interaction, rather than a
token-generation step. The alignment error is $e_t=r-y_t$, and the
objective is to reduce this error within a finite interaction budget
using prompt interventions alone.

A candidate intervention must be represented in the predictor to make
its effect distinguishable from that of other candidates. To formalize
this interface, we assume a command-conditioned prompting protocol
\begin{equation}
    u_t=\mathcal{P}(x,\mathcal{H}_t,c_t),
    \qquad
    c_t\in\mathcal{C}\subseteq\mathbb{R},
    \label{eq:prompt_protocol}
\end{equation}
where $\mathcal{P}$ constructs a prompt from a numerical command $c_t$
and the current dialogue. The desired output $r$ and the applied
command $c_t$ have distinct roles: the former defines the alignment
objective, whereas the latter is selected by the controller. In
ordinary target-following trajectories, $c_t$ equals the requested
target. Predictive control allows the command to differ from $r$ to
compensate for the response predicted by the learned model. This
formulation concerns the specified prompting protocol, not arbitrary
textual interventions with an unspecified encoding.

\subsection{Finite-Memory Koopman Dynamics}
\label{sec:memory_koopman}

The current attribute alone need not contain sufficient information to
predict the effect of another prompt. For example, the same sentence
length can occur while successive revisions are becoming longer or
while they are becoming shorter. We therefore represent the interaction
using a finite window of attribute observations,
\begin{equation}
    s_t=
    \begin{bmatrix}
        y_t & y_{t-1} & \cdots & y_{t-L}
    \end{bmatrix}^{\top}
    \in\mathbb{R}^{L+1},
    \qquad
    y_t=C s_t,
    \label{eq:memory_state}
\end{equation}
where $L\geq 1$ is the number of previous observations retained and
$C=[1,0,\ldots,0]$ selects the current attribute. This representation
exposes recent levels and trends to the predictor without requiring
access to hidden activations. It is an approximation to the information
needed for prediction, not an exact reconstruction of the full dialogue
state.

Under this finite-memory approximation, define the command-conditioned
Koopman operator acting on an observable $f$ by
\begin{equation}
    (\mathcal{K}_{c}f)(s)
    =
    \mathbb{E}\!\left[
        f(s_{t+1})
        \mid s_t=s,\,c_t=c
    \right].
    \label{eq:koopman_operator}
\end{equation}
The operator is linear in $f$, although the underlying prompt--response
dynamics can be nonlinear. We approximate its action in a learned,
finite-dimensional observable space using an encoder $E_{\phi}$ and
a decoder $D_{\psi}$:
\begin{equation}
    \begin{aligned}
        \xi_t
        &= E_{\phi}(s_t),
        & \xi_t &\in\mathbb{R}^{d_{\xi}},\\
        \widehat{\xi}_{t+1}
        &= K\xi_t+B c_t+b,\\
        \widehat{s}_{t+1}
        &= D_{\psi}(\widehat{\xi}_{t+1}),
        & \widehat{y}_{t+1}
        &= C\widehat{s}_{t+1}.
    \end{aligned}
    \label{eq:koopman_model}
\end{equation}
Here, $K\in\mathbb{R}^{d_{\xi}\times d_{\xi}}$ describes latent
evolution, $B\in\mathbb{R}^{d_{\xi}\times 1}$ describes command
dependence, and $b\in\mathbb{R}^{d_{\xi}}$ is a bias. We write
$g_{\psi}=C\circ D_{\psi}$ for the attribute readout.

The nonlinear encoder and decoder allow the observed dynamics to be
nonlinear despite the affine latent transition. The shared $K$ and
affine command dependence are modeling restrictions, rather than
consequences of the operator's linearity. In particular, we do not
assume an exact finite-dimensional Koopman representation of the LLM.
Moreover, $g_{\psi}$ need not be linear, so linear latent dynamics do
not imply a linear map from latent states to output attributes.

\subsection{Learning the Koopman Predictor}
\label{sec:learning}

Given recorded prompt--response trajectories, we construct transition
tuples $\mathcal{D}=\{(s_t,c_t,s_{t+1})\}$ using overlapping history
windows. Each command label is the command actually applied before the
successor observation. For trajectories collected with a fixed
requested target, that target supplies the command label throughout
the trajectory. Let
\begin{equation}
    q_t
    =
    K E_{\phi}(s_t)+B c_t+b.
    \label{eq:training_transition}
\end{equation}
We jointly optimize state reconstruction, latent consistency, and
next-state prediction:
\begin{equation}
    \begin{aligned}
        \mathcal{L}_{\mathrm{rec}}
        &=
        \mathbb{E}_{\mathcal{D}}\!\left[
            \left\|
                s_t-D_{\psi}(E_{\phi}(s_t))
            \right\|_2^2
        \right],\\
        \mathcal{L}_{\mathrm{lin}}
        &=
        \mathbb{E}_{\mathcal{D}}\!\left[
            \left\|
                E_{\phi}(s_{t+1})-q_t
            \right\|_2^2
        \right],\\
        \mathcal{L}_{\mathrm{pred}}
        &=
        \mathbb{E}_{\mathcal{D}}\!\left[
            \left\|
                s_{t+1}-D_{\psi}(q_t)
            \right\|_2^2
        \right],\\
        \mathcal{L}
        &=
        \lambda_{\mathrm{rec}}\mathcal{L}_{\mathrm{rec}}
        +
        \lambda_{\mathrm{lin}}\mathcal{L}_{\mathrm{lin}}
        +
        \lambda_{\mathrm{pred}}\mathcal{L}_{\mathrm{pred}}.
    \end{aligned}
    \label{eq:training_objective}
\end{equation}
The reconstruction term preserves the information in the memory state.
The latent consistency term encourages the encoded successor to agree
with the affine transition. The prediction term ties this transition
to accuracy in the observed state space. These objectives are
complementary: latent consistency alone admits uninformative constant
encodings, whereas reconstruction alone does not require the
representation to support prediction. Training updates
$(\phi,\psi,K,B,b)$, while $\theta$ remains unchanged. The learned
predictor is held fixed during test-time control.

\subsection{Predictive Test-Time Control}
\label{sec:control}

At each iteration, let $\mathcal{C}_t\subseteq\mathcal{C}$ be a
nonempty finite set of admissible candidate commands. We evaluate each
candidate by holding its command fixed over a prediction horizon of
$H$ interactions under the same prompting protocol, truncated to the
remaining interaction budget. Starting from the current observed
memory,
\begin{equation}
    \begin{aligned}
        \widehat{\xi}_{t\mid t}(c)
        &= E_{\phi}(s_t),\\
        \widehat{\xi}_{t+h\mid t}(c)
        &=
        K\widehat{\xi}_{t+h-1\mid t}(c)+Bc+b,\\
        \widehat{y}_{t+h\mid t}(c)
        &=
        g_{\psi}\!\left(
            \widehat{\xi}_{t+h\mid t}(c)
        \right),
        \qquad h=1,\ldots,H.
    \end{aligned}
    \label{eq:rollout}
\end{equation}
The latent trajectory also has the explicit form
\begin{equation}
    \widehat{\xi}_{t+h\mid t}(c)
    =
    K^h E_{\phi}(s_t)
    +
    \sum_{j=0}^{h-1}K^j(Bc+b).
    \label{eq:closed_form_rollout}
\end{equation}
Consequently, candidate evaluation requires latent propagation and
attribute decoding rather than additional generations from the
target LLM.

We score each candidate using its predicted trajectory's discrepancy
from the desired attribute:
\begin{equation}
    \begin{aligned}
        \widehat{J}_t(c)
        &=
        \sum_{h=1}^{H}w_h
        \left|
            \widehat{y}_{t+h\mid t}(c)-r
        \right|^2,\\
        c_t^{\star}
        &\in
        \operatorname*{arg\,min}_{c\in\mathcal{C}_t}
        \widehat{J}_t(c),
        \qquad w_h\geq 0.
    \end{aligned}
    \label{eq:selection}
\end{equation}
The controller executes only the next intervention
$u_t=\mathcal{P}(x,\mathcal{H}_t,c_t^{\star})$. After receiving
$o_{t+1}$, it evaluates $y_{t+1}=V(o_{t+1})$, shifts the memory window,
and repeats the optimization. Thus, the constant-command assumption
applies only within a hypothetical rollout; the executed command can
change at every real interaction.

Control begins once an initialization prefix provides $L+1$
observations and terminates when the measured error meets a specified
tolerance or the interaction budget is exhausted. Each new prediction
is initialized from actual observations rather than a previously
predicted memory. Candidate simulation adds no target-LLM queries,
although any cost of constructing candidates or evaluating responses
must still be included in the overall inference budget. Accurate
prediction on fixed-command training trajectories alone does not
establish accuracy after command changes; control additionally depends
on coverage of the command--history pairs encountered at inference
time.

\subsection{Prediction Error and Intervention Selection}
\label{sec:analysis}

We give a conditional analysis connecting predictive accuracy to the
quality of candidate selection. For a candidate $c$, let $s_h^c$ and
$y_h^c$ denote the actual memory state and attribute after $h$ further
interactions that repeatedly apply $c$ from the current state. Define
$\xi_h^c=E_{\phi}(s_h^c)$, with $s_0^c=s_t$. Suppose that, uniformly
over the candidate rollouts,
\begin{equation}
    \begin{aligned}
        \xi_{h+1}^c
        &=K\xi_h^c+Bc+b+\eta_h^c,
        &
        \|\eta_h^c\|_2
        &\leq\varepsilon_{\mathrm{dyn}},\\
        y_h^c
        &=g_{\psi}(\xi_h^c)+\nu_h^c,
        &
        |\nu_h^c|
        &\leq\varepsilon_{\mathrm{rec}}.
    \end{aligned}
    \label{eq:residual_assumptions}
\end{equation}
These residuals include imperfect finite-memory closure, model
approximation, and generation variability. Assume further that
$\|K\|_2\leq\kappa$ and $g_{\psi}$ is $L_g$-Lipschitz on the relevant
latent region. For stochastic generation, the following bounds apply
on the event that the stated residual bounds hold.

\paragraph{Multi-step prediction error.}

Because the predicted and actual latent rollouts start from the same
encoding, their initial discrepancy is zero. Subtracting their
recursions yields
$\|d_{h+1}\|_2
\leq\kappa\|d_h\|_2+\varepsilon_{\mathrm{dyn}}$,
where
$d_h=\xi_h^c-\widehat{\xi}_{t+h\mid t}(c)$.
Iterating this inequality and applying the Lipschitz readout gives
\begin{equation}
    \left|
        y_h^c-\widehat{y}_{t+h\mid t}(c)
    \right|
    \leq
    \varepsilon_{\mathrm{rec}}
    +
    L_g\varepsilon_{\mathrm{dyn}}
    \sum_{j=0}^{h-1}\kappa^j
    \;\triangleq\;\Delta_h.
    \label{eq:prediction_bound}
\end{equation}
This bound exhibits the trade-off in choosing $H$: a longer horizon
accounts for more distant consequences, but requires accurate dynamics
over a longer rollout. If $\kappa<1$, the geometric sum is bounded by
$1/(1-\kappa)$. No such contraction is assumed for the learned model.

\paragraph{Reliability of candidate selection.}

Define the actual held-command cost
$J_t(c)=\sum_{h=1}^{H}w_h|y_h^c-r|^2$.
Suppose both actual and predicted attribute errors are bounded by $R$
for all candidates and prediction steps. The difference-of-squares
identity and Eq.~\eqref{eq:prediction_bound} imply
\begin{equation}
    \left|
        J_t(c)-\widehat{J}_t(c)
    \right|
    \leq
    2R\sum_{h=1}^{H}w_h\Delta_h
    \;\triangleq\;\delta_J.
    \label{eq:cost_bound}
\end{equation}
Let
$c_t^{\mathrm{opt}}
\in\operatorname*{arg\,min}_{c\in\mathcal{C}_t}J_t(c)$.
Using the uniform cost bound twice and the definition of
$c_t^{\star}$ yields
\begin{equation}
    \begin{aligned}
        J_t(c_t^{\star})
        &\leq
        \widehat{J}_t(c_t^{\star})+\delta_J\\
        &\leq
        \widehat{J}_t(c_t^{\mathrm{opt}})+\delta_J\\
        &\leq
        \min_{c\in\mathcal{C}_t}J_t(c)+2\delta_J.
    \end{aligned}
    \label{eq:selection_bound}
\end{equation}
Thus, uniformly accurate predictions support near-optimal selection
within the evaluated candidate set. The statement concerns the current
held-command comparison, not the cumulative cost of the subsequently
replanned closed loop. It does not guarantee that the candidate set
contains an intervention capable of reaching $r$, nor that the closed
loop converges. Re-observation resets the rollout's initial latent
error; it does not eliminate model mismatch or guarantee exact
target tracking.